\chapter*{คำนำ}
\addcontentsline{toc}{chapter}{คำนำ}

รายวิชา \textbf{4011111 ฟิสิกส์สำหรับจุลชีววิทยา} (\textbf{Physics for Microbiology}) จำนวน 3(2-3-4) หน่วยกิต เป็นรายวิชาบังคับในกลุ่มวิชาแกนวิทยาศาสตร์ของหลักสูตรวิทยาศาสตรบัณฑิต สาขาวิชาจุลชีววิทยา คณะวิทยาศาสตร์และเทคโนโลยี มหาวิทยาลัยราชภัฏรำไพพรรณี โดยมีเป้าหมายสำคัญเพื่อสร้างรากฐานความเข้าใจด้านปรากฏการณ์ทางกายภาพและกลไกการทำงานของเครื่องมือวิทยาศาสตร์ที่นักจุลชีววิทยาต้องใช้งานจริงในห้องปฏิบัติการ ตั้งแต่การปิเปตสารละลายระดับไมโครลิตร การเคลื่อนที่ของแบคทีเรียในอาหารเลี้ยงเชื้อเหลว การแยกเซลล์ด้วยเครื่องปั่นเหวี่ยง (\textbf{Centrifuge}) การฆ่าเชื้อด้วยหม้อนึ่งความดันไอ (\textbf{Autoclave}) การย้อมสีแกรมและการแยกดีเอ็นเอด้วยสนามไฟฟ้า (\textbf{Gel Electrophoresis}) การเพิ่มอำนาจแยกชัดของกล้องจุลทรรศน์ด้วยน้ำมันช่วยส่อง (\textbf{Immersion Oil}) ไปจนถึงการควบคุมจุลินทรีย์ด้วยรังสีอัลตราไวโอเลตและการวัดความหนาแน่นเซลล์ด้วยเครื่องสเปกโตรโฟโตมิเตอร์ (\textbf{Spectrophotometer})

ปัญหาสำคัญที่พบในการจัดการเรียนการสอนวิชาฟิสิกส์ให้แก่นักศึกษาสาขาจุลชีววิทยา คือ นักศึกษาจำนวนมากมีความกังวลต่อสมการคณิตศาสตร์ที่ซับซ้อนและแคลคูลัส ดังนั้น ผู้เขียนจึงได้ออกแบบและเรียบเรียงเอกสารคำสอนและคู่มือปฏิบัติการเล่มนี้ขึ้นใหม่ทั้งหมดตามแนวคิด \textbf{ฟิสิกส์ฉบับเข้าใจง่ายไร้แคลคูลัส} (\textbf{Low-Math \& High-Concept Physics}) โดยตัดการใช้แคลคูลัสและตรีโกณมิติที่ซับซ้อนออกทั้งหมด แล้วแทนที่ด้วย \textbf{เทคนิคสามเหลี่ยมสูตรคำนวณ} การเปรียบเทียบสัดส่วนตรงและสัดส่วนผกผัน และการอ่านกราฟเชิงกายภาพควบคู่ไปกับตัวอย่างโจทย์ในบริบทห้องปฏิบัติการจุลชีววิทยาที่มีการแสดงวิธีทำทีละขั้นตอนอย่างชัดเจน

นอกจากนี้ เอกสารคำสอนเล่มนี้ยังได้รับการออกแบบตามกรอบมาตรฐานการจัดการศึกษาที่มุ่งเน้นผลลัพธ์การเรียนรู้ (\textbf{Outcome-Based Education}) โดยเชื่อมโยงผลลัพธ์การเรียนรู้ของหลักสูตรสาขาวิชาจุลชีววิทยา (\textbf{PLO1 ถึง PLO6}) เข้ากับผลลัพธ์การเรียนรู้ของรายวิชา (\textbf{CLO1 ถึง CLO4}) อย่างเคร่งครัด ครอบคลุมแผนการสอนรายสัปดาห์ตลอด 15 สัปดาห์ (บรรยาย 2 ชั่วโมง และปฏิบัติการ 3 ชั่วโมง รวม 5 ชั่วโมงต่อสัปดาห์) พร้อมใบงานปฏิบัติการทดลองทั้ง 13 การทดลองอยู่ในเล่มเดียวกัน เพื่อให้นักศึกษาสามารถใช้เรียนรู้ทั้งภาคทฤษฎีและภาคปฏิบัติได้อย่างสมบูรณ์

ผู้เขียนหวังเป็นอย่างยิ่งว่า เอกสารคำสอนและคู่มือปฏิบัติการรายวิชา 4011111 ฟิสิกส์สำหรับจุลชีววิทยา เล่มนี้ จะช่วยให้นักศึกษาสาขาวิชาจุลชีววิทยาเกิดความมั่นใจ มีเจตคติที่ดีต่อวิชาฟิสิกส์ และสามารถนำหลักการทางฟิสิกส์ไปประยุกต์ใช้ในการทำงานวิจัยและการปฏิบัติงานในห้องปฏิบัติการจุลชีววิทยาได้อย่างมืออาชีพ

\vspace{1.2cm}
\begin{flushright}
    \textbf{ผู้ช่วยศาสตราจารย์ ดร.ชีวะ ทัศนา}\\[0.2cm]
    สาขาวิชาฟิสิกส์ คณะวิทยาศาสตร์และเทคโนโลยี\\
    มหาวิทยาลัยราชภัฏรำไพพรรณี\\
    ภาคเรียนที่ 2 ปีการศึกษา 2569
\end{flushright}
